% =====================================================================
%  MODELO EM BRANCO -- seção de implementação em formato metodológico
%
%  Use uma seção para cada "percurso" completo de implementação (por
%  exemplo: abastecimento de água; esgotamento sanitário; estações de
%  monitoramento). A seção descreve o percurso DO COMEÇO AO FIM, na ordem
%  em que as coisas acontecem, e não atividade por atividade.
%
%  Como usar:
%   1. Copie este arquivo para capitulos/<subcomponente>/0X-<nome>.tex
%   2. Substitua tudo o que estiver entre << >>
%   3. Inclua o arquivo no capitulo.tex do subcomponente com \input{...}
%
%  Regras de redação:
%   - TEXTO CORRIDO. Parágrafos de 5 a 10 linhas, que expliquem o que
%     será feito, por quem, em que ordem, com que produto e por quê.
%     Evite listas de tópicos com frases curtas.
%   - Poucos títulos: \section para o percurso e \subsection para cada
%     fase. Não use \subsubsection nem \paragraph.
%   - Um fluxograma por fase relevante (ver pasta figuras/):
%       * sequência ....... figuras/fluxograma-captacao-projeto.tex
%       * raias (quem faz o quê) .. figuras/fluxograma-obras.tex
%       * linha do tempo com produtos .. figuras/fluxograma-modelo-gestao.tex
%   - Datas de 2026-2027: use o POA e remeta ao Anexo A.
%   - Cite o código da atividade do MOP no texto (\atividade{X.X.X.X}).
%   - Quadros longos de etapas vão para o Apêndice B.
%   - Inconsistências: \pendencia{...} logo após o parágrafo afetado.
% =====================================================================
\section{Implementação de <<percurso: ex. sistemas de abastecimento de água>>}
\label{sec:<<rotulo>>}

<<Parágrafo de abertura: em quantas fases o percurso se divide, do que
trata cada uma e quais atividades do MOP ele reúne (\atividade{X.X.X.X}).
Remeter à figura geral do percurso.>>

\begin{figure}[htbp]
  \centering
  \caption{<<Implementação de ...>>, do começo ao fim}
  \label{fig:<<rotulo>>}
  \fluxograma{<<fluxograma-geral-do-percurso>>}
  \fonte{Elaborado pelo IAT (2026) com base em \citeonline{mop2026}.}
\end{figure}

% ---------------------------------------------------------------------
\subsection{<<Fase 1: ex. Seleção e adesão>>}
\label{subsec:<<rotulo-fase1>>}

<<Como a fase começa (o que a antecede), quem conduz, o que será feito em
cada passo, com que critérios ou normas, que produto resulta e como ele é
comprovado. Incluir as datas do POA quando houver.>>

<<Segundo parágrafo, se necessário: decisões, pontos de verificação,
riscos específicos e o que a fase entrega para a fase seguinte.>>

\begin{figure}[htbp]
  \centering
  \caption{<<Título do fluxograma da fase>>}
  \label{fig:<<rotulo-fase1>>}
  \fluxograma{<<arquivo>>}
  \fonte{Elaborado pelo IAT (2026) com base em \citeonline{mop2026}.}
\end{figure}

% ---------------------------------------------------------------------
\subsection{<<Fase 2>>}
\label{subsec:<<rotulo-fase2>>}

<<...>>

% ---------------------------------------------------------------------
% Fechamento do percurso: quadro-síntese curto (uma linha por atividade)
% ---------------------------------------------------------------------
O \autoref{qua:sintese-<<rotulo>>} resume, para cada atividade, o
responsável, o prazo, o produto e a forma de comprovação.

\begin{quadro}[htbp]
\caption{Síntese das atividades de <<percurso>>}
\label{qua:sintese-<<rotulo>>}
\footnotesize
\renewcommand{\arraystretch}{1.25}
\begin{tabularx}{\textwidth}{|P{2.9cm}|P{2.0cm}|Q{1.3cm}|L|L|P{1.4cm}|}
\hline
\cab{Atividade} & \cab{Responsável} & \cab{Prazo} & \cab{Produto} &
\cab{Evidência} & \cab{Validação} \\ \hline
<<X.X.X.X Nome>> & <<>> & <<Anos X a Y>> & <<>> & <<>> & <<>> \\ \hline
\end{tabularx}
\fonte{Adaptado de \citeonline{mop2026}.}
\end{quadro}
